\newpage
\section{Related Work}

\paragraph{Formal Verification of AI Systems.}
There is a growing body of work on applying formal methods to AI, including verification of neural network robustness~\cite{katz2017reluplex,wang2023alphaverus}, safe reinforcement learning~\cite{fulton2018safe}, and formal specification of AI goals~\cite{russell2019human}. Our work differs in focusing on the \textit{foundational epistemic layer} of a knowledge-genesis system, rather than on learned model behavior.

\paragraph{Multi-Prover and Cross-Verification.}
The idea of using multiple provers to increase trust is well-established in high-assurance software (e.g., seL4's dual verification in Isabelle/HOL and Coq~\cite{klein2009sel4}). Our contribution is the application of this principle to the novel domain of AI foundational ontologies, and the specific integration of SymPy, Z3, Lean~4, and Coq for the task of atomicity verification.

\paragraph{Certified Programming and Proof Assistants.}
Projects like the DeepSpec consortium~\cite{deepspec} aim to build fully verified software stacks. The \Apeiron{} Framework's layered design is inspired by such compositional verification approaches, though applied to a cognitive rather than a traditional software architecture.

\paragraph{AI Alignment via Formal Constraints.}
Recent work on ``provably beneficial AI''~\cite{everitt2021provably} proposes using formal methods to constrain AI behavior. Our work complements this by providing a mechanism to verify the \textit{substrate} on which such constraints operate.